\subsection{Solution to Problem 10: Complex Numbers with Argument Condition}

\begin{solution}
We are given $\arg(z/w) = -\pi/2$, so $z/w$ is a negative purely imaginary number. Write
\[
\frac{z}{w} = -ik \qquad \text{for some real } k > 0,
\]
giving $z = -ikw$. Then
\begin{align*}
\left|\frac{z-w}{z+w}\right|
&= \left|\frac{-ikw - w}{-ikw + w}\right|
= \left|\frac{-w(1+ik)}{w(1-ik)}\right| \\
&= \left|\frac{1+ik}{1-ik}\right|
= \frac{|1+ik|}{|1-ik|}
= \frac{\sqrt{1+k^2}}{\sqrt{1+k^2}}
= 1.
\end{align*}

Alternatively, since $z/w$ is purely imaginary, $z$ and $w$ are perpendicular in the Argand plane. The points $z-w$ and $z+w$ are diagonals of a rectangle with sides parallel to the real and imaginary directions, so they have equal length and the ratio is $1$.

Therefore $\left|\dfrac{z-w}{z+w}\right| = 1$.
\end{solution}

\begin{takeaways}
\begin{itemize}
    \item $\arg(z/w) = -\pi/2$ means $z/w$ is negative purely imaginary
    \item Write $z = -ikw$ with real $k > 0$ and simplify the modulus
    \item Geometric view: numerator and denominator are equal diagonals of a rectangle
\end{itemize}
\end{takeaways}

\begin{remark}
An easy solution is geometric: the numerator and denominator are diagonals of a rectangle.
\end{remark}
